\documentclass[11pt]{article}
\usepackage[T1]{fontenc}
\usepackage{amsmath,amssymb,amsthm}
\usepackage[margin=1in]{geometry}
\usepackage[hidelinks]{hyperref}

\title{Literature-Implied Answer: Convex Weak-Star Derived Sets Above \texorpdfstring{$\omega+1$}{omega+1}}
\author{Run fa\_banach\_001}
\date{2026-06-19}

\begin{document}
\maketitle

\noindent\textbf{Status.}
\texttt{literature\_implied\_answer (source Question 2; source Question 1 remains open).}
This is a literature-status packet, not an original proof from the run.

\section*{Source Question}

The source paper is Zdenek Silber, \emph{Weak$^*$ derived sets of convex sets in duals of non-reflexive spaces}, arXiv:2101.10684, J. Funct. Anal. 281 (2021), no. 12, Paper No. 109259 \cite{Sil21}.  In Section 3, ``Remarks and open problems'', Silber asks two questions.  The second asks:
\[
\text{if } X \text{ is non-reflexive and quasi-reflexive, are there convex subsets of } X^*
\text{ with order higher than } \omega+1?
\]

\section*{Supporting Theorem}

The supporting paper is Mikhail I. Ostrovskii, \emph{Weak$^*$ closures and derived sets for convex sets in dual Banach spaces}, arXiv:2112.04670 \cite{Ost21}.  The abstract and Theorem~1 state the successor-ordinal result for convex subsets of duals: for every non-reflexive Banach space and every countable successor ordinal, there is a convex subset of the dual whose weak-star derived-set process reaches the weak-star closure precisely at that successor stage.

\section*{Identification}

Ostrovskii's theorem applies to every non-reflexive Banach space, hence in particular to every non-reflexive quasi-reflexive Banach space.  Choosing any countable successor ordinal strictly larger than $\omega+1$, for instance $\omega+2$, gives a convex subset of $X^*$ with order higher than $\omega+1$.  Therefore Silber's second question has a positive answer in the later literature.

The supporting paper cites Silber's 2021 paper and presents its theorem as a further development of the convex-set weak-star derived-set theory.  The answer to the quasi-reflexive question is therefore immediate from the theorem, but the supporting paper does not appear to restate that specific question verbatim; this is why the packet is placed under \texttt{literature\_implied\_answers}.

\section*{Limitations}

This packet does not answer Silber's first question asking whether the order of a convex subset of the dual of a separable Banach space can be a limit ordinal.  Ostrovskii's paper explicitly records that this limit-order question remains open in its final comment section.

\section*{Search and Provenance}

The cheap run indexes were searched for arXiv:2101.10684, weak derived sets, derived sets, convex sets, and related phrases.  No prior packet or attempt was found.  Local source search and web search then identified arXiv:2112.04670 as the decisive later theorem.

\begin{thebibliography}{9}
\bibitem{Sil21}
Z. Silber,
\emph{Weak$^*$ derived sets of convex sets in duals of non-reflexive spaces},
J. Funct. Anal. 281 (2021), no. 12, Paper No. 109259; arXiv:2101.10684.

\bibitem{Ost21}
M. I. Ostrovskii,
\emph{Weak$^*$ closures and derived sets for convex sets in dual Banach spaces},
arXiv:2112.04670.
\end{thebibliography}

\end{document}
